\begin{afterword}
\summaryhead

The post opens with satire. The aliens of film and television look like humans in costume, and the post mocks the idea that convergent evolution explains this: Earth's own animals look less like us the longer ago their lineages split from ours. The plain explanation is that costumed humans are the easy way out. The real failure, the author says, is in mind: aliens who think like English speakers of the present day.

We understand other people, the author argues, by ``empathic inference.'' Our ancestors had to model only minds like their own, and a human brain is too complex to model except by running a shadow of its states on one's own. This works because humans share one psychology, with the same emotions and facial expressions everywhere. It fails for minds with other emotions, with no sense of humor, or with no emotions at all, and natural selection is such a real alien. Anthropomorphism cannot be dropped by an act of will. The author praises Jack Vance and \textsc{The Mote in God's Eye} for writing real aliens, recalls a Star Trek episode whose aliens had rewritten the preamble to the U.S. Constitution, and concludes that the inability to imagine the alien is the inability to see one's own humanity as special.

\responsehead

The satire works, and its target is real: screen aliens often think like the people who wrote them. The step from shape to mind, and the \textsc{Psycho} anecdote about one's own culture as a special case, make the thesis plain. The problems are in the psychology offered to explain the failure.

The account of how we understand other minds is one side of a debate, stated as fact. That we ``evolved to predict Other Minds by putting ourselves in their shoes'' is what philosophers and psychologists call simulation theory; its rival holds that we use a tacit theory of mind. The post does not mention the rival. Its argument, that other brains are too complex to model ``neuron-by-neuron,'' rules out one alternative to simulation and not the other: an abstract model of goals and beliefs. The previous post in the book, ``Belief in Intelligence,'' uses such a model to predict where a friend will drive.

The evidence for human sameness is stated more firmly than the research allows. That the six emotions are shown ``by the same facial expressions'' in every known culture was disputed in 1994 on grounds of method, and a large review in 2019 found that expressions vary substantially across cultures. The argument from sexual reproduction, that complex adaptations must be shared, does not depend on the facial evidence.

The case for minds without emotions is only half made. The imagined objection is that ``any intelligent mind powerful enough to build complex machines'' must have something like emotions. The answer, natural selection, shows that building complex machines does not require emotions. It does not show that an intelligent mind can lack them, since the author elsewhere calls evolution ``stupid in an absolute sense'' and does not count it as an intelligence. The paragraph on humor has a related gap. The claim that evolution must give intelligent minds a sense of humor is raised and not answered; a weaker one is answered with a comparison to shouting English abroad.

In short: a sharp satire of screen aliens, whose account of how we understand other minds takes one side of an open debate as settled, and whose counterexample to minds needing emotions is, by the author's own account, not a mind.
\end{afterword}
